% -*- mode: latex; coding: utf-8 -*-
%-----------------------------------------------------------------------------
% pxORL -- User manual and API reference
% Part of P3CTeX. Author: Scvria.
%
% Build from directory tex/ (with TEXINPUTS=./latex;./code;./doc;):
%   pdflatex -interaction=nonstopmode -output-directory=doc doc/pxORL.tex
%   pdflatex -interaction=nonstopmode -output-directory=doc doc/pxORL.tex
%-----------------------------------------------------------------------------
\documentclass[11pt,a4paper]{article}
\PassOptionsToPackage{colorlinks,urlcolor=blue,linkcolor=blue}{hyperref}
\usepackage[T1]{fontenc}
\usepackage[utf8]{inputenc}
\usepackage[margin=22mm]{geometry}
\usepackage{lmodern}
\usepackage[catalan]{babel}
\usepackage{pxORL}
\usepackage{hyperref}

\title{El paquet \texttt{pxORL}\\
  \large Generaci\'o d'orles acad\`emiques i taulers de graduaci\'o per a \LaTeX2e\ i P3CTeX}
\author{Scvria\\P3CTeX / UOC PEC Repository}
\date{2026/10/27 -- Versi\'o 0.1}

\begin{document}
\maketitle

\begin{abstract}
  \texttt{pxORL} proporciona un sistema complet i automatitzat per generar orles acad\`emiques i
  taulers de graduaci\'o basats en llistes d'estudiants i fotografies, aprofitant la identitat visual,
  els colors i la tipografia corporativa de la UOC i del paquet \textbf{P3CTeX}. Ofereix c\`alcul
  autom\`atic de columnes i dimensions, suport per a marcadors de posici\'o (\emph{placeholders})
  quan manquen imatges, i diversos presets predefinits.
\end{abstract}

\tableofcontents
\newpage

\section{C\`arrega i activaci\'o}

Dins d'un document amb la classe \texttt{P3CTeX}, activeu el m\`odul \texttt{ORL} a les opcions de classe:

\begin{verbatim}
\documentclass[ORL,PAC,default]{P3CTeX}
\end{verbatim}

Tamb\'e es pot utilitzar en qualsevol document \LaTeX\ independent, habitualment configurat en orientaci\'o apa\"isada (\emph{landscape}):

\begin{verbatim}
\usepackage[a4paper,landscape]{geometry}
\usepackage{pxORL}
\end{verbatim}

\section{Comanda principal: \texttt{\textbackslash pxOrla} i \texttt{\textbackslash pxORLprint}}

La generaci\'o de l'orla es realitza mitjan\c{c}ant la instrucci\'o \texttt{\textbackslash pxOrla} (o la seva variant t\`ecnica \texttt{\textbackslash pxORLprint}):

\begin{verbatim}
\pxOrla[<opcions locals>]{
  {Nom Estudiant 1}/fitxer_foto1,
  {Nom Estudiant 2}/fitxer_foto2,
  {Nom Estudiant 3}/none
}
\end{verbatim}

\begin{itemize}
  \item Cada entrada s'especifica en el format \texttt{\{Nom\}/foto}.
  \item Si no es disposa de la fotografia d'un estudiant, utilitzeu \texttt{none} o una ruta inexistent: \texttt{pxORL} dibuixar\`a autom\`aticament un marc silueta professional sense aturar la compilaci\'o.
  \item Les opcions locals permeten ajustar el t\'itol, el subt\'itol, el peu, les columnes i els colors.
\end{itemize}

\section{Configuraci\'o i presets}

\subsection{\texttt{\textbackslash pxORLsetup}}

Permet establir la configuraci\'o persistent que s'aplicar\`a a totes les orles posteriors:

\begin{verbatim}
\pxORLsetup{
  title       = {Promoció 2022--2026},
  subtitle    = {Grau en Enginyeria Informàtica},
  footer      = {Universitat Oberta de Catalunya},
  columns     = 5,
  photos-path = {fotos/}
}
\end{verbatim}

\subsection{Presets predefinits}

\texttt{pxORL} incorpora quatre configuracions d'estil:

\begin{description}
  \item[\texttt{p3c-uoc-orla}:] Estil corporatiu de la UOC amb blaus primari i secundari, orles de doble l\'inia i tipografia formal. \'Es el valor per defecte.
  \item[\texttt{p3c-classic}:] Disseny cl\`assic amb marcs ornamentats suaus i disposici\'o equilibrada.
  \item[\texttt{p3c-compact}:] Disseny compacte d'alta densitat, idoni per a promocions nombroses amb molts alumnes per p\`agina.
  \item[\texttt{p3c-formal}:] Disseny d'alt contrast amb tipografia serifa per a actes acad\`emics de cloenda solemne.
\end{description}

Per activar un preset:

\begin{verbatim}
\pxORLusepreset{p3c-compact}
\end{verbatim}

Tamb\'e podeu registrar nous presets amb \texttt{\textbackslash pxORLsavepreset\{<nom>\}\{<claus>\}}.

\end{document}
